\section{Target and exact event representation}
For VUS, let $y\in\{0,1\}^n$ contain both classes and $s\in\mathbb{R}^n$ contain finite
scores after binary64 conversion. Larger scores indicate greater anomaly
evidence. Let $P_0=\sum_i y_i$, and let the $K$ maximal original anomaly
intervals have inclusive endpoints $[a_k,b_k]$ in temporal order. Widths are
the integers $w=0,\ldots,L$. The continuous label $a_i^{(w)}$ is one on
original positives. An interval contributes $\sqrt{1-d/w}$ outside its
endpoints at integer distance $1\le d\le\lfloor w/2\rfloor$; contributions
are summed and clipped at one. At $w=0$ no ramp is added. Expanded intervals
that overlap merge, while adjacent nonoverlapping intervals remain separate.
This is the original dynamic VUS target pinned to the official
\texttt{RangeAUC\_volume\_opt} source at revision
\texttt{467bace0}, configured for exhaustive thresholds~\cite{vusreference}.
It is not shared by all software named VUS.

For a threshold $t$, define
\begin{align}
 p_i(t)&=\ind\{s_i\ge t\}, & C(t)&=\sum_i p_i(t),\\
 A(t)&=\sum_i y_i p_i(t), &
 B(t)&=\sum_i(1-y_i)a_i^{(w)}p_i(t),\\
 T(t)&=A(t)+B(t), & P(t)&=P_0+B(t)/2,\\
 N(t)&=n-P(t).
\end{align}
Let $H(t)$ count expanded components hit by at least one predicted positive,
and let $K_w$ count those components. The curve coordinates are
\begin{align}
 R(t)&=\min\{T(t)/P(t),1\}\,H(t)/K_w,\\
 F(t)&=(C(t)-T(t))/N(t), & \pi(t)&=T(t)/C(t).
\end{align}
Precision is used only at nonempty retained thresholds, where $C(t)>0$;
the initial origin requires no precision value.
The denominator is positive: $P(t)\le(n+P_0)/2<n$. ROC integrates trapezoids
from the origin through complete distinct descending score thresholds to
the explicit terminal $(1,1)$. PR uses the right-step sum
$\sum_j(R_j-R_{j-1})\pi_j$, with $R_0=0$. Each VUS scalar is the arithmetic
mean over every width. Signed horizontal ROC increments are preserved; the
derivation does not replace a dynamic denominator by a fixed one.

Let $S_w=\{i:a_i^{(w)}>0\}$, $Q_w=|S_w|$, and $Q=Q_L$. For an active event
with score $v_j$, precompute the complete global counts
$c_j^>=|\{i:s_i>v_j\}|$ and $c_j^\ge=|\{i:s_i\ge v_j\}|$. They include
positions outside $S_w$ and outside $S_L$.

\begin{lemma}[Inactive threshold runs]
Between consecutive active score events at a fixed width, $T,B,P,N,H,R$
are constant. Omitted thresholds contribute a horizontal ROC rectangle and
zero PR area. At an active score, its complete global numerical tie is a
single transition.
\end{lemma}
\begin{proof}
A newly admitted inactive point has zero weight and cannot hit an expanded
component. Only $C$ changes. The horizontal increments
$(\Delta C)/N$ telescope at fixed height $R$. All positions tied with an
active event enter at its global threshold, so both counts must retain them.
\end{proof}

If the previous retained state is $(T_o,N_o,R_o,c_o)$, event $j$ contributes
\begin{align}
 f_j^-&=(c_j^>-T_o)/N_o,\\
 \Delta\mathrm{ROC}_j&=
 (c_j^>-c_o)R_o/N_o+(F_j-f_j^-)(R_j+R_o)/2,\\
 \Delta\mathrm{PR}_j&=(R_j-R_o)T_j/c_j^\ge.
\end{align}
The initial state is $(0,n-P_0,0,0)$, and afterwards $c_o=c_j^\ge$.
Connecting retained endpoints without the intervening rectangle is
incorrect. For example, $y=(1,0,1,0)$, $s=(3,2,1,0)$ and $L=0$ give ROC
area $0.625$; connecting support endpoints alone gives $0.8125$.

\begin{lemma}[Two-distance geometry]
The weight of a non-original point is determined by its distances $d_1,d_2$
to the nearest two distinct original intervals. It is zero if
$d_1>\lfloor w/2\rfloor$, one if $d_2\le\lfloor w/2\rfloor$, and
$\sqrt{1-d_1/w}$ otherwise.
\end{lemma}
\begin{proof}
Every active ramp has weight at least $1/\sqrt2$. Two contributions therefore
clip to one. If exactly one interval contributes, its distance determines
the weight. Distances to ordered intervals are monotone moving away on each
side, so two neighbors on each side contain the nearest two intervals.
\end{proof}

\begin{theorem}[Support-event evaluation]
\label{thm:support-events}
In a unit-cost model with score comparisons and metric arithmetic, including
square roots, both exhaustive scalar integrals can be computed in
$O(n\log(Q+1)+(L+1)Q)$ time and $O(n+Q)$ auxiliary words.
\end{theorem}
\begin{proof}
Extract original intervals and maximum support in $O(n)$ time. Sort the $Q$
support scores and group equal values. Binary-search all $n$ scores among
these groups to obtain complete global counts $c^>,c^\ge$ without sorting
off-support scores, at cost $O(n\log(Q+1))$. The ordered-interval geometry
can be prepared within this bound. At each width, compute support weights,
component membership, grouped masses and first component-hit events in
$O(Q+K)$ time. A scan of the at most $Q$ event groups integrates the
preceding rectangles, complete tied transitions and terminal ROC tail.
The inactive-run lemma makes this identical to the exhaustive target in
exact metric arithmetic. Stream the width mean using input-, support- and
interval-sized arrays. Since $K\le Q\le n$, the stated bounds follow.
\end{proof}

\begin{proposition}[Fixed-binary64 preparation]
\label{prop:radix}
For finite binary64 scores, complete global tie counts, the score ordering
of maximum support, and its two-distance geometry can be prepared in
$O(n+Q+K)$ time and $O(n+Q)$ auxiliary words. Discrete work uses a word-RAM
with input indices, counts, widths, and intermediate index arithmetic
fitting words. Metric arithmetic, including square roots, has unit cost.
\end{proposition}
\begin{proof}
Normalize signed zeros. Complement a negative score's bit pattern and flip
the sign bit of a nonnegative score, then complement the result to obtain a
descending numerical key. Eight stable byte-counting passes order these
64-bit keys in $O(n)$ time. Scan complete ties and retain their starting and
ending global ranks. Filter the permutation to maximum-support positions.
A temporal scan with a monotone interval pointer computes nearest distances
and owners in $O(Q+K)$ time. All arrays are input- or support-sized.
\end{proof}
Radix sorting and floating-point key transforms are established
techniques~\cite{demaine2011radix,herf2001radix}. This assumption does not give
linear sorting of arbitrary real scores or a bit-complexity bound for growing
precision. Conversion defines numerical ties, including collisions between
distinct large integer inputs.

The same support-event width kernel combined with this preparation gives
$O(n+(L+1)Q)$ time. This is the basic Sparse VUS evaluator; its improvement
over a dense rank sweep depends on $Q$, rather than recall saturation.
The released PR-only implementation retains $L+1$ scalar width results,
adding $O(L)$ words; streaming removes this storage. Input-sized ranking
arrays remain necessary in this construction, so the memory is not
support-only.
